\chapter{Configuration Reference}
\label{app:config_reference}
The system behavior is largely controlled by TOML and YAML files. Table \ref{tab:config_reference_app} highlights key parameters.

\begin{figure}[H]
    \centering
\begin{lstlisting}[style=yamlStyle]
# Example sequencer.yaml configuration
server:
  host: "0.0.0.0"
  port: 3000
batching:
  max_batch_size: 100
  timeout_interval_ms: 2000
scheduling:
  policy_type: "FCFS"
  timeboost_window_ms: 5000
\end{lstlisting}
    \caption{Example Configuration Dump for Sequencer}
\end{figure}

\begin{table}[H]
\centering
\caption{Key Configuration Parameters}
\label{tab:config_reference_app}
\small
\begin{tabularx}{\textwidth}{>{\hsize=1.0\hsize}X >{\hsize=1.0\hsize}X >{\hsize=0.5\hsize}X >{\hsize=1.5\hsize}X}
\toprule
\textbf{Parameter} & \textbf{Location} & \textbf{Default} & \textbf{Effect} \\
\midrule
\code{max\_batch\_size} & \texttt{sequencer.toml} & 100 & Triggers a batch seal when the tx pool reaches this size. \\
\addlinespace
\code{timeout\_interval\_ms} & \texttt{sequencer.toml} & 5000 & Triggers a batch seal if the interval expires, regardless of size. \\
\addlinespace
\code{policy\_type} & \texttt{sequencer.toml} & FCFS & Determines the transaction ordering algorithm (e.g., FCFS or FeePriority). \\
\addlinespace
\code{da\_mode} & \texttt{submitter.yaml} & calldata & Swaps the L1 data posting mechanism between standard calldata and EIP-4844 blobs. \\
\bottomrule
\end{tabularx}
\end{table}